\hypertarget{native-282-nexus6}{}\label{native-282-nexus6}
\Article{6}
Bevor wir die Folgerungen hieraus allgemein untersuchen, erschien es uns zweckmäßig, einige besondere Fälle auszuarbeiten, damit sich die Art der Argumentation leichter verstehen lässt. Greifen wir also zunächst die Form $xx+yy$ wieder auf, für die die Anzahl der Darstellungen einer unbestimmten Zahl den mittleren Wert $=\pi$ erhält. Die Anzahl der tatsächlichen Darstellungen einer gegebenen Zahl lässt sich ohne Schwierigkeit aus den in den Untersuchungen über höhere Arithmetik aufgestellten allgemeinen Prinzipien bestimmen. Bezeichnen wir mit $fA$ die Anzahl der Darstellungen der Zahl $A$, so ist diese $=4$, wenn $A=1$ oder $2$ oder eine Potenz von zwei ist; $=8$, wenn $A$ eine Primzahl der Form $4n+1$ oder das Produkt
\SourceLeaf{283}
einer solchen Primzahl mit einer Potenz von zwei ist; $=0$, wenn $A$ eine Primzahl der Form $4n+3$ oder durch eine solche Primzahl, jedoch nicht durch deren Quadrat teilbar ist; schließlich \emph{allgemein}
\[
\begin{aligned}
\text{entweder}\quad&=4(\alpha+1)(\beta+1)(\gamma+1)\ldots\\
\text{oder}\quad&=0
\end{aligned}
\]
Dabei wird die Zahl $A$ auf die Form $2^\mu Sa^\alpha b^\beta c^\gamma\ldots$ gebracht; $a,b,c$ usw. bezeichnen verschiedene Primzahlen der Form $4n+1$ und $S$ das Produkt der Primzahlen der Form $4n+3$, die ein- oder mehrmals unter den Faktoren der Zahl $A$ vorkommen, falls es solche gibt. Es gilt der erste oder zweite Wert, je nachdem, ob die Zahl $S$ ein Quadrat oder kein Quadrat ist. Somit hängt $fA$ offenbar allein davon ab, wie die Primzahlen $3,5,7,11,13$ usw. unter den Faktoren der Zahl $A$ vorkommen, sodass allgemein
\[
fA=4(3)\pdot(5)\pdot(7)\pdot(11)\pdot(13)\dotsfour
\]
zu setzen ist, wenn wir annehmen, dass die Werte der Zeichen $(3),(5),(7)$ usw. so gewählt werden, dass für eine Primzahl $p$ gilt:

\begin{description}
\item[erstens] $(p)=1$, wenn $p$ die Zahl $A$ nicht teilt
\item[zweitens] $(p)=\alpha+1$, wenn $p$ die Form $4n+1$ hat und $p^\alpha$ die höchste Potenz ist, die $A$ teilt
\item[drittens] $(p)=0$, wenn $p$ die Form $4n+3$ hat und der Exponent der höchsten Potenz von $p$, die $A$ teilt, ungerade ist; schließlich
\item[viertens] $(p)=1$, wenn $p$ die Form $4n+3$ hat und der Exponent der höchsten Potenz von $p$, die $A$ teilt, gerade ist.
\end{description}
Offenbar ist der erste Fall im zweiten und vierten enthalten.

Auf diese Weise verlaufen die Glieder der Folge $f1,f2,f3,f4$ usw. sehr unregelmäßig, obwohl sich der mittlere Wert $=\pi$ umso genauer ergeben muss, je mehr Glieder genommen werden. Die Summe $f1+f2+f3+\ldots+fA$ bezeichnen wir mit $FA$.

\Article{7}
Setzen wir nun allgemein $fm+f3m=f'm$, so erkennt man leicht, dass
\[
f'A=4(5)\pdot(7)\pdot(11)\pdot(13)\ldots
\]
\SourceLeaf{284}
gilt, das heißt, $f'A$ wird von der Beziehung von $A$ zum Teiler $3$ unabhängig sein, sodass die Unregelmäßigkeit der Folge $f'1,f'2,f'3,f'4,f'5,f'6$ usw. sowohl später beginnt als auch viel geringer ist. Setzen wir ferner
\[
f'1+f'2+f'3+f'4+\text{usw.}+f'm=F'm
\]
so gilt
\[
\begin{aligned}
F'3A&=F3A+f3+f6+f9+\ldots+f9A\\
&=F3A+FA
\end{aligned}
\]
Daraus folgt leicht, dass für ins Unendliche wachsendes $A$
\[
F'3A=4\pi A
\]
\begingroup\emergencystretch=1em\displaywidowpenalty=0
zu setzen ist, oder dass der mittlere Wert der Glieder der Folge\hfil\break
\mbox{$f'1,f'2,f'3,f'4$ usw.}
\[
=\tfrac43\pi
\]
ist.
\par\endgroup

Ebenso ergibt sich, wenn wir allgemein $-f'm+f'5m=f''m$ setzen,
\[
f''A=4(7)(11)(13)\ldots
\]
das heißt, aus der neuen Folge $f''1,f''2$ usw. verschwinden die Schwankungen, die von der Beziehung zur Zahl $5$ abhängen. Setzen wir die Summe
\[
f''1+f''2+f''3+\ldots+f''m=F''m
\]
so gilt
\[
F''5m=-F'm+F'5m
\]
woraus folgt, dass für ins Unendliche wachsendes $m$
\[
F''5m=\tfrac43\pi\pdot4m
\]
zu setzen ist, oder dass der mittlere Wert der Glieder der Folge $=\tfrac43\pdot\tfrac45\pi$ ist.

Gehen wir in derselben Weise weiter, indem wir neue Folgen bilden und dabei nacheinander die Faktoren $(7),(11),(13),(17)$ usw. entfernen, so nähern sich diese immer mehr der Unveränderlichkeit, und ihre mittleren Werte erhalten nacheinander die neuen Faktoren $\tfrac87,\tfrac{12}{11},\tfrac{12}{13},\tfrac{16}{17},\tfrac{20}{19}$ usw. Die Nenner sind dabei die Primzahlen in ihrer natürlichen Reihenfolge, die Zähler dagegen um eins größer oder kleiner, je nachdem, ob jene die Form $4n-1$ oder ${}^{.}4n+1$ haben. Da nun bei Fortsetzung dieses Verfahrens ins Unendliche der konstante Wert
\SourceLeaf{285}
$4$ dem mittleren Wert immer näher kommen muss, erhalten wir
\[
4=\pi\pdot\tfrac43\pdot\tfrac45\pdot\tfrac87\pdot\tfrac{12}{11}\pdot\tfrac{12}{13}\ldots\text{ bis ins Unendliche}
\]
oder
\[
\pi=4\pdot\frac3{3+1}\pdot\frac5{5-1}\pdot\frac7{7+1}\pdot\frac{11}{11+1}\pdot\frac{13}{13-1}\ldots
\]
Entwickelt man die einzelnen Brüche in unendliche Reihen,
\begingroup\postdisplaypenalty=10000
\[
\begin{aligned}
\frac3{3+1}&=1-\frac13+\frac19-\frac1{27}+\ldots\\
\frac5{5-1}&=1+\frac15+\frac1{25}+\frac1{125}+\ldots\\
\frac7{7+1}&=1-\frac17+\frac1{49}-\frac1{343}+\ldots\\
&\text{usw.}
\end{aligned}
\]
so lässt sich das Produkt leicht in
\endgroup
\[
1-\tfrac13+\tfrac15-\tfrac17+\tfrac19-\tfrac1{11}+\tfrac1{13}\dotsfour
\]
entwickeln; die Summe dieser Reihe ist bekanntlich $=\tfrac14\pi$. Tatsächlich hatte der berühmte \textsc{Euler} schon vor langer Zeit auf umgekehrtem Weg hieraus die Gleichheit zwischen $\tfrac14\pi$ und dem unendlichen Produkt $\tfrac43\pdot\tfrac45\pdot\tfrac87\ldotp\ldotp$ hergeleitet (\emph{Introd. in analys. inf.}, Bd. I, Kap. xv, Artikel~285).

\clearpage
\Article{8}
Betrachten wir an zweiter Stelle die Form $xx+2yy$, für die die Anzahl der Darstellungen einer unbestimmten Zahl den mittleren Wert $=\frac\pi{\surd2}$ haben wird. Bezeichnen wir mit $fA$ die Anzahl der Darstellungen einer gegebenen Zahl $A$ durch diese Form, so ist diese $=2$ für $A=1$ oder $A=2$ oder immer dann, wenn $A$ eine Potenz von zwei ist; ferner $fA=4$, wenn $A$ eine der Primzahlen ist, deren quadratischer Rest $-2$ ist, also eine der Formen $8k+1,8k+3$, etwa $A=3,11,17,19,41,43$ usw.; schließlich $fA=0$, wenn $A$ eine Primzahl ist, deren quadratischer Nichtrest $-2$ ist, etwa aus der Folge $5,7,13,23,29,31$ usw., also entweder der Form $8k+5$ oder der Form $8k+7$. Allgemein ist jedoch
\[
\begin{aligned}
\text{entweder}\quad fA&=2(\alpha+1)(\beta+1)(\gamma+1)\ldots\\
\text{oder}\quad fA&=0
\end{aligned}
\]
zu setzen. Dabei wird die Zahl $A$ auf die Form $2^\mu Sa^\alpha b^\beta c^\gamma\ldots$ gebracht; $a,b,c$ usw. bezeichnen verschiedene Primzahlen der Formen $8k+1,8k+3$, dagegen $S$ das Produkt der
\SourceLeaf{286}
übrigen Zahlen (der Formen $8k+5,8k+7$), falls solche unter den Faktoren der Zahl $A$ vorkommen. Es gilt der erste oder zweite Wert, je nachdem, ob $S$ ein Quadrat oder kein Quadrat ist. Durch ganz ähnliche Schlussfolgerungen wie im vorigen Artikel gehen wir nun von der Folge $f1,f2,f3,f4,f5$ usw., also $2,2,4,2,0,2$ usw., nacheinander zu anderen über, die immer länger \emph{konstant} bleiben. Ihre \emph{mittleren} Werte sind nacheinander $\frac\pi{\surd2}$, $\frac\pi{\surd2}\pdot\frac23$, $\frac\pi{\surd2}\pdot\frac23\pdot\frac65$, $\frac\pi{\surd2}\pdot\frac23\pdot\frac65\pdot\frac87$ usw.; wir gehen dabei so vor, dass wir zu der Gleichung
\[
2=\frac\pi{\surd2}\pdot\frac23\pdot\frac65\pdot\frac87\pdot\frac{10}{11}\pdot\frac{14}{13}\pdot\frac{16}{17}\pdot\frac{18}{19}\ldots
\]
gelangen, in der die Nenner die natürliche Folge der Primzahlen bilden, die Zähler dagegen um eins kleiner sind, wenn die Nenner die Form $8k+1$ oder $8k+3$ haben, und um eins größer, wenn die Nenner die Form $8k+5$ oder $8k+7$ haben.
\BookRule{.30\textwidth}
